\chapter{Summary and Outlook}

This thesis focused on the advanced 2D and 3D characterisation of clinker phases and hydrated cementitious binders by combining a variety of modern imaging and analytical techniques.
The main focus was the application of modern scanning electron microscopy (\ac{SEM}) imaging techniques, which were combined with various complementary methods.
The primary goal was to improve the understanding of cement hydration and microstructural development by establishing and refining robust evaluation methods and workflows.

\begin{description}

    \item[\ac{ArBIB} preparation technique:]
    The importance of advanced specimen preparation techniques like \acl{ArBIB} (\acs{ArBIB}) polishing for obtaining artefact-free surfaces for high-resolution \ac{SEM} imaging is especially important for ever more complex microstructures in modern binder systems.
    It was demonstrated that \ac{ArBIB} sectioning preserves the native pore structure without the need for resin embedding.
    Furthermore, it produces smooth surfaces even at the interfaces between soft and hard materials.
    For example, it enables the application of mechanical testing methods such as \acl{NI} (\acs{NI}) within an \acl{ITZ} (\acs{ITZ}) on unembedded hydrated concrete samples.

    \item[Mechanical phase properties:] %Correlative microscopy of clinker phases
    Correlating \ac{NI}, \acl{EDX} (\acs{EDX}), and \acl{BSE} (\acs{BSE}) enabled to determine phase-specific mechanical properties and highlighted differences between unhydrated clinker and hydration products near an \ac{ITZ}.
    These results can be used to investigate and compare the development of \ac{ITZ} structures of upcoming composite cements and \textsc{Portland} cement.

    \item[Chemical composition of clinker phases:]
    A correlative workflow combining \ac{SEM}-\ac{EDX} and \acl{LAICPMS} was established to map major, minor, and trace elements in unhydrated clinker.
    This approach allowed precise localisation of trace elements within specific clinker phases, demonstrating, for example, enrichment of \ce{Ba}, \ce{V} and \ce{Rb} in \ac{C2S} and of \ce{Mn} in the interstitial phase.
    Combining multiple techniques proved a powerful strategy to overcome individual-method limitations in resolution and sensitivity.

    \item[Advanced \ac{FIBnt} workflows:]
    A refined lift-out technique for \acl{FIBnt} (\acs{FIBnt}) was optimised for cementitious materials to minimise common artefacts (charging, redeposition and shadowing), enabling acquisition of high-quality 3D datasets.
    A post-processing workflow was developed to register, denoise and segment correlative \ac{BSE} and \ac{EDX} 3D datasets, utilising 3D superpixel clustering for phase identification.
    Laser preparation was also introduced to access larger volumes and targeted regions of interest such as the \ac{ITZ}.
    This will be useful to obtain an insight into complex nano and microstructures of modern binder systems while imaging a statistical more representative volume.

    \item[\ac{HLT} measurements:]
    A tool was developed to automatically measure the \ac{HLT} from 2D \ac{SEM}-\ac{BSE} images of polished sections of hydrated clinkers.
    The method was applied to \ac{C3S}, \ac{C2S} and their mixtures, revealing the influence of phase composition and age on hydration kinetics.
    Pixel classification-based segmentation was shown to be more reliable than simple thresholding.
    A statistical analysis defined the minimum representative area required for phase quantification, which depends on sample heterogeneity.

    \item[Particle distance measurements:]
    Cryo-\acs{FIB}-\ac{SEM} was used to analyse inter-par\-ti\-cle distances in fresh cement pastes.
    A modified version of the \ac{HLT} measurement algorithm was used to quantify particle distances.
    The method objectively demonstrated the dispersing effects of \acl{PCE} and \acl{PUS} (\acs{PUS}) on the microstructure of cement pastes by quantifying shifts in the particle-distance distributions.

    \item[\ac{EBSD}-\ac{HLT} correlation:]
    Correlation of automated \ac{HLT} measurements with \ac{EBSD}-determined crystal orientations of \ac{C3S} revealed no significant preference for specific crystallographic directions in \ac{CSH} growth.
    Instead, \ac{HLT} development appears to be primarily governed by the available pore space.
    A correlation between \ac{EBSD} measurement density and crystal orientation was, however, observed.
    The cause of reduced measurement density for certain orientations remains unclear and further measurements are required.
    Nevertheless, this finding demonstrates an additional application of \ac{EBSD}.

    \item[Comparative pore structure analysis:]
    A comparative study of pore analysis methods (2D and 3D-\ac{SEM}, \acl{MIP}, and \acl{LTDSC} (\acs{LTDSC})) was conducted on reference materials (porous glass) and hydrated cement pastes.
    \ac{LTDSC} proved effective for analysing nanoscale gel porosity, showing an increase in gel-pore volume with hydration time for both \ac{C3S} and \ac{C2S}.
    While \ac{C2S} exhibited a higher overall capillary porosity, \ac{C3S} showed a substantially higher gel-pore volume.
    High-resolution \ac{SEM} and \ac{FIBnt} imaging of \ac{ArBIB}-prepared cross-sections allowed direct observation of nano-porosity within inner \ac{CSH}, yielding pore-size distributions comparable to indirect methods.
    However, the inherent limitations of the individual techniques make direct comparisons challenging.

\end{description}

These results demonstrate that combining various imaging and analytical techniques yields significantly more comprehensive information than individual methods by mitigating their respective limitations and leveraging their specific strengths.
Crucially, the reliability of these analyses depends heavily on specimen preparation techniques, such as \ac{ArBIB} milling and cryo-fixation, which are essential for preserving the native microstructure.
Furthermore, the application of modern image and data processing algorithms not only accelerates the image evaluation but also offers insights extending well beyond basic image segmentation and phase quantification.
This transition from a qualitative observation to a quantitative statistical analysis, as exemplified by the automated \ac{HLT} measurements and particle distance analyses, enables a more objective and robust characterisation of cementitious materials.

With annual worldwide cement production of \qty{3.8}{\giga\tonne} \cite{usgs.2026}, even mar\-gi\-nal improvements in binder efficiency, durability, and sustainability can have a massive environmental and economic impact.
To achieve these improvements, the historical `lack of knowledge' mentioned by \textsc{Anderegg} and \textsc{Hubbell} \cite{anderegg.1929} regarding cement hydration must be continually addressed through modern analytical capabilities.
By providing robust, quantitative tools to decode these increasingly complex microstructures, the methods and workflows developed in this thesis lay the groundwork for future research:
\begin{itemize}
    \item The automated \ac{HLT} measurement can be applied to a wider range of binder systems, such as \textsc{Portland} and composite cements. These results can be used to validate hydration models, such as the level-set method. Furthermore, the \ac{HLT} measurement can be combined with other techniques, as demonstrated with \ac{EBSD}, to obtain deeper insights into the factors influencing the hydration process.
    \item The correlative 3D datasets from \ac{FIBnt} provide a basis for microstructural modelling, simulation of transport properties, and validation of 2D-to-3D conversion algorithms.
    \item Since \ac{FIBnt} volumes prepared solely by \ac{FIB} are inherently small, it was demonstrated that the preparation of larger volumes can be significantly accelerated by utilising the \acl{fs-laser} technique.
    \item The cryo-\ac{FIB}-\ac{SEM} technique for analysing fresh pastes can be extended to study the very early stages of hydration and the influence of different admixtures or physical preparation steps like \ac{PUS} on particle interactions in situ.
    \item The most challenging and time-consuming step remains the image segmentation task.
    Further development of machine learning-based segmentation for both 2D and 3D datasets will be crucial to increase the throughput and objectivity of microstructural analyses.
    To achieve this goal, the cement community needs to increase its efforts to freely publish annotated datasets.
\end{itemize}